\chapter{Conclusion and Future Work}

\section*{Conclusion}
FL-IDS shows that the five parts of the system work together on real IoT/IIoT traffic: a boosting classifier, a federated autoencoder backstop, a cosine-similarity trust filter, Flower-based federated learning, and SDN mitigation. Each part has a measured contribution. All planned phases are complete (Table~\ref{tab:workplan}); what remains is the research paper.

\begin{itemize}
  \item \textbf{Accuracy target met.} With service-port, payload-content and flow features and square-root class weights, the LightGBM classifier reaches \textbf{99.7\% accuracy and 0.983--0.985 macro F1} when trained on 1.43 million packets and tested on 397 thousand packets from connections that never appear in training (a flow-grouped 80/20 split of all 1.99 million unique packets). That is up from 96.4\% and 0.889 with the original features. The result holds across three seeds, and it holds when the classifier is trained on only the 5\% server calibration set (99.7\% / 0.979).
  \item \textbf{The whole federated system improves with the new features.} A real run with ten client processes reaches 0.999 attack macro-recall (was 0.889), and the cascade detects never-seen attack classes at 0.943 macro-recall (was 0.736), against 0.801 for the classifier alone. The backstop's value is now clearest on unseen attacks: on known attacks the classifier alone has the higher detection F1 (0.998 vs.\ 0.961).
  \item \textbf{The trust filter keeps federated training usable under poisoning.} Attack detection stays at 0.997--1.000 and zero-day detection at 0.90--0.94 up to 40\% malicious clients, while plain FedAvg's autoencoder diverges at 30\%. The filter fails at 50\%, as expected for median-based defences.
  \item \textbf{Detection drives enforcement.} Phase B derives the same features live (99.8\% of live rows match training rows exactly); a replayed Backdoor capture was classified with 0.997 confidence and produced a drop rule, confirmed from the switch's flow table.
  \item \textbf{Data quality.} The released Edge-IIoTset CSV contains a placeholder label leak (one feature with an AUROC of 1.0), first reported by Galal~\cite{galal2026provenance}, and column-shifted MITM rows, which we found and repaired. Results trained directly on it are likely inflated.
\end{itemize}

\section*{Limitations}
The service-port feature benefits from the testbed's lack of benign HTTP traffic. Two attack classes are detected less well as zero-day attacks with the new features (MITM, and Uploading for the classifier alone), and attacker survival at 40\% rose from 24\% to 36\% in a single-seed run. Flow features are computed per 120-second flow record, so classification is not per-packet. Mitigation in Phase B is capture-level rather than inline. Two attack models were tested against the trust filter (sign-flip, and a white-box stealth attacker that the filter mostly misses), and the stealth, pruning and clipping experiments use one seed. At full scale the false-alarm rate is 2.5\% to 16.3\% of normal packets but only 0.34\% to 1.04\% of normal connections, because a few very long connections dominate the packet count; the trust filter and the autoencoder backstop are weaker than at 60{,}000 rows; MITM, a single connection, is scored only by the zero-day experiment and the live replay; and two of the 14 captures are classified as the wrong type in the live replay (both are still blocked). The federated system now trains on all usable rows (Section~\ref{sec:fullpool}), but its bootstrap classifier starts from the 5\% server set; the standalone classifier trained on 1.43 million rows (Table~\ref{tab:datascale}). The classifier update ships a small number of confirmed alert rows to the server, which relaxes strict data locality.

\section*{Future Work}
\begin{enumerate}
  \item Make mitigation streaming: classify flows incrementally (for example after the first $n$ packets) so a rule can be installed while an attack is still running.
  \item Defend against the white-box stealth attacker (Section~\ref{sec:followups}), which the filter mostly misses, for example with a small trusted root dataset as in FLTrust~\cite{cao2021fltrust}; measure zero-day detection under it; and test other adaptive strategies and attackers above 50\%.
  \item Add more seeds (now three) and confidence intervals to the ablation, poisoning sweep and zero-day experiments, and repeat the zero-day columns inside the first two at the extra seeds.
  \item Find out why the trust filter excludes fewer attackers at full scale, and repeat the stealth-attacker, pruning and clipping experiments at full scale over more seeds.
  \item Validate on a second dataset that contains benign web traffic, to measure how much the service-port feature depends on the testbed.
  \item Reduce the per-round communication as well: weight pruning (Objective 3, Section~\ref{sec:followups}) shrinks the model, with autoencoder recall 0.744 dense and 0.688 at 70\% sparsity, but clients still send dense updates. Sparse updates, and zero-day detection of pruned models, remain to be done.
  \item Federate the classifier itself (for example with histogram-based federated boosting) so the calibration set no longer needs to be held centrally.
  \item Deploy on physical edge hardware and a hardware OpenFlow switch.
\end{enumerate}
